\documentclass{article}
\usepackage{amsmath,amssymb,amsthm}
\usepackage{algorithm}
\usepackage{algpseudocode}
\usepackage{fullpage}
\usepackage{hyperref}
\newtheorem{theorem}{Theorem}
\newtheorem{lemma}{Lemma}
\newtheorem{assumption}{Assumption}
\newtheorem{remark}{Remark}
\begin{document}

\section*{Adaptive Variance‑Reduced Gradient (AVRG)}

\section*{Mathematical Formulation}
Let $f:\mathbb{R}^d\rightarrow\mathbb{R}$ be the empirical risk
\[
f(x)=\frac{1}{n}\sum_{i=1}^{n}f_i(x),\qquad 
g_i(x)=\nabla f_i(x).
\]
At iteration $t$ a minibatch $\mathcal{B}_t\subset\{1,\dots,n\}$ of size $b$ is sampled uniformly at random and the stochastic gradient
\[
\tilde g_t:=\frac{1}{b}\sum_{i\in\mathcal{B}_t} g_i(x_t)
\]
is computed.  

\paragraph{Recursive variance estimator}
\[
v_t=(1-\alpha)v_{t-1}+\alpha\,\tilde g_t,\qquad \alpha\in(0,1].
\tag{1}
\]
The estimator is biased because of the exponential moving average.  A bias‑corrected version is
\[
\hat v_t:=\frac{v_t}{1-(1-\alpha)^t}.
\tag{2}
\]

\paragraph{Adaptive stepsize}
\[
\eta_t=\frac{\eta}{1+\beta\|\hat v_t\|},\qquad \eta>0,\;\beta\ge 0.
\tag{3}
\]

\paragraph{Parameter update}
\[
x_{t+1}=x_t-\eta_t \,\tilde g_t .
\tag{4}
\]

The algorithm is initialized with $x_0\in\mathbb{R}^d$, $v_{-1}=0$, and runs for $T$ iterations or until a stopping criterion is met.

\section*{Pseudocode}
\begin{algorithm}[H]
\caption{Adaptive Variance‑Reduced Gradient (AVRG)}
\begin{algorithmic}[1]
\State \textbf{Input:} stepsize $\eta>0$, momentum weight $\alpha\in(0,1]$, adaptivity $\beta\ge0$, max iter $T$.
\State Initialize $x_0$, $v_{-1}=0$.
\For{$t=0,\dots,T-1$}
    \State Sample minibatch $\mathcal{B}_t$ of size $b$ uniformly.
    \State $\tilde g_t\gets\frac{1}{b}\sum_{i\in\mathcal{B}_t}\nabla f_i(x_t)$.
    \State $v_t\gets (1-\alpha)v_{t-1}+ \alpha \tilde g_t$.\Comment{Recursive estimator}
    \State $\hat v_t\gets v_t/(1-(1-\alpha)^{t+1})$.\Comment{Bias correction}
    \State $\eta_t\gets \dfrac{\eta}{1+\beta\|\hat v_t\|}$.
    \State $x_{t+1}\gets x_t-\eta_t \tilde g_t$.
\EndFor
\State \textbf{Output:} $x_T$ (or the averaged iterate $\bar x_T$).
\end{algorithmic}
\end{algorithm}

\section*{Dry Run Example}
Consider the one‑dimensional quadratic
\[
f(w)=(w-3)^2,\qquad \nabla f(w)=2(w-3).
\]
Set $w_0=0$, $\eta=0.1$, $\alpha=0.5$, $\beta=0$, and use a full‑batch gradient ($b=n$).  

\begin{center}
\begin{tabular}{c|c|c|c|c}
$t$ & $\tilde g_t= \nabla f(w_t)$ & $v_t$ (by (1)) & $\hat v_t$ (by (2)) & $w_{t+1}=w_t-\eta\tilde g_t$\\ \hline
0 & $2(0-3)=-6$ & $v_0=0.5(-6)=-3$ & $\hat v_0=-3/ (1-(0.5)^1) =-6$ & $w_1=0-0.1(-6)=0.6$\\
1 & $2(0.6-3)=-4.8$ & $v_1=0.5(-3)+0.5(-4.8)=-3.9$ & $\hat v_1=-3.9/(1-(0.5)^2)=-5.2$ & $w_2=0.6-0.1(-4.8)=1.08$\\
2 & $2(1.08-3)=-3.84$ & $v_2=0.5(-3.9)+0.5(-3.84)=-3.87$ & $\hat v_2=-3.87/(1-(0.5)^3)=-4.86$ & $w_3=1.08-0.1(-3.84)=1.464$\\
\end{tabular}
\end{center}
The objective values decrease: $f(w_0)=9$, $f(w_1)=5.76$, $f(w_2)=3.686$, $f(w_3)=2.336$.

\newpage

\section*{Assumptions}
\begin{assumption}[Smoothness]\label{ass:smooth}
Each component $f_i$ is $L$‑smooth:
\[
\|\nabla f_i(x)-\nabla f_i(y)\|\le L\|x-y\|,\qquad\forall x,y\in\mathbb{R}^d .
\]
Consequently $f$ is $L$‑smooth.
\end{assumption}

\begin{assumption}[Polyak–Łojasiewicz (PL) condition]\label{ass:pl}
There exists $\mu>0$ such that for all $x$,
\[
\frac{1}{2}\|\nabla f(x)\|^{2}\ge \mu\bigl(f(x)-f^\star\bigr),
\]
where $f^\star=\min_x f(x)$.
\end{assumption}

\begin{assumption}[Unbiased stochastic gradients]\label{ass:unbiased}
\[
\mathbb{E}_{\mathcal{B}_t}\bigl[\tilde g_t\mid x_t\bigr]=\nabla f(x_t),\qquad
\mathbb{E}_{\mathcal{B}_t}\bigl[\|\tilde g_t-\nabla f(x_t)\|^{2}\mid x_t\bigr]\le\sigma^{2}.
\]
\end{assumption}

\begin{assumption}[Bound on the recursive estimator]\label{ass:bound}
The moving‑average weight satisfies $\alpha\in(0,1]$ and we define
\[
c_\alpha:=\frac{1}{1-(1-\alpha)}=\frac{1}{\alpha}.
\]
Hence $\|\hat v_t\|\le c_\alpha\max_{0\le s\le t}\|\tilde g_s\|$ almost surely.
\end{assumption}

\section*{Main Convergence Theorem}
\begin{theorem}[Linear convergence under PL]\label{thm:linear}
Assume \ref{ass:smooth}–\ref{ass:bound} and choose a constant base stepsize $\eta$ such that
\[
0<\eta\le \min\Bigl\{\frac{1}{L},\;\frac{1}{2\beta c_\alpha G}\Bigr\},
\]
where $G:=\sup_{x}\|\nabla f(x)\|<\infty$ (finite for quadratic or bounded domains).  
Then the iterates generated by AVRG satisfy
\[
\mathbb{E}\bigl[f(x_T)-f^\star\bigr]\le
\bigl(1-\tfrac{\mu\eta}{2}\bigr)^{T}\bigl[f(x_0)-f^\star\bigr]
\;+\;
\frac{\eta\sigma^{2}}{2\mu}.
\tag{5}
\]
In particular, if $\sigma=0$ (full‑batch or variance‑reduced case) the algorithm converges linearly with rate $1-\mu\eta/2$.
\end{theorem}

\section*{Theoretical Properties}
\begin{itemize}
\item \textbf{Stability.} The adaptive stepsize (3) is bounded below by $\eta/(1+\beta c_\alpha G)$ and above by $\eta$, guaranteeing that the recursion never uses a stepsize larger than the standard constant stepsize.
\item \textbf{Hyperparameter sensitivity.}
    \begin{itemize}
        \item Larger $\alpha$ gives a faster response of $v_t$ to recent gradients, decreasing $c_\alpha$ and allowing a larger admissible $\beta$.
        \item $\beta=0$ reduces AVRG to standard SGD with momentum‑style EMA but without adaptivity.
    \end{itemize}
\item \textbf{Comparison to baselines.}
    \begin{itemize}
        \item When $\beta=0$, Theorem~\ref{thm:linear} recovers the classic linear rate of SGD under PL (cf. \cite{karimi2016linear}).
        \item With $\beta>0$, the stepsize automatically shrinks in high‑variance regimes, yielding the same asymptotic rate while providing better practical robustness.
    \end{itemize}
\end{itemize}

\section*{Discussion}
The theorem shows that the EMA‑based variance estimator together with bias correction does not deteriorate the PL‑induced linear convergence.  The extra term $\eta\sigma^{2}/(2\mu)$ mirrors the variance floor typical for stochastic methods; employing a SARAH/SPIDER inner loop to compute $\tilde g_t$ makes $\sigma^{2}=O(\epsilon^{2})$ and therefore the error floor can be driven arbitrarily low.

\newpage

\section*{Preliminaries (Definitions and Lemmas)}
\begin{lemma}[Descent Lemma for $L$‑smooth functions]\label{lem:descent}
For any $x,y\in\mathbb{R}^{d}$,
\[
f(y)\le f(x)+\langle\nabla f(x),y-x\rangle+\frac{L}{2}\|y-x\|^{2}.
\]
\end{lemma}
\begin{proof}
Standard; follows from integrating the Lipschitz condition on the gradient. \hfill $\square$
\end{proof}

\begin{lemma}[Bound on adaptive stepsize]\label{lem:stepbound}
Under Assumption~\ref{ass:bound} and the choice of $\eta$ in Theorem~\ref{thm:linear},
\[
\frac{\eta}{1+\beta c_\alpha G}\le \eta_t\le \eta,\qquad\forall t\ge0.
\]
\end{lemma}
\begin{proof}
From (3) and $\|\hat v_t\|\le c_\alpha G$ we have 
$1+\beta\|\hat v_t\|\le 1+\beta c_\alpha G$, hence $\eta_t\ge \eta/(1+\beta c_\alpha G)$. 
Since $\|\hat v_t\|\ge0$, $\eta_t\le\eta$. \hfill $\square$
\end{proof}

\section*{Main Proof (Step‑by‑Step Derivation)}
We prove Theorem~\ref{thm:linear}.  Throughout we condition on the sigma‑algebra $\mathcal{F}_t:=\sigma(x_0,\dots,x_t)$.

\paragraph{[Step 1] Apply the descent lemma with $y=x_{t+1}$ and $x=x_t$.}
\begin{align}
f(x_{t+1})
&\le f(x_t)+\langle\nabla f(x_t),x_{t+1}-x_t\rangle+\frac{L}{2}\|x_{t+1}-x_t\|^{2}
\notag\\
&= f(x_t)-\eta_t\langle\nabla f(x_t),\tilde g_t\rangle+\frac{L}{2}\eta_t^{2}\|\tilde g_t\|^{2}.
\tag{6}
\end{align}

\paragraph{[Step 2] Take expectation w.r.t.\ $\tilde g_t$ conditional on $\mathcal{F}_t$.}
Using unbiasedness (Assumption~\ref{ass:unbiased}),
\[
\mathbb{E}\bigl[\langle\nabla f(x_t),\tilde g_t\rangle\mid\mathcal{F}_t\big]
= \|\nabla f(x_t)\|^{2}.
\]
Similarly,
\[
\mathbb{E}\bigl[\|\tilde g_t\|^{2}\mid\mathcal{F}_t\bigr]
= \|\nabla f(x_t)\|^{2}+ \mathbb{E}\bigl[\|\tilde g_t-\nabla f(x_t)\|^{2}\mid\mathcal{F}_t\bigr]
\le \|\nabla f(x_t)\|^{2}+\sigma^{2}.
\tag{7}
\]

\paragraph{[Step 3] Insert (7) into (6) and use Lemma~\ref{lem:stepbound}.}
\begin{align}
\mathbb{E}\!\bigl[f(x_{t+1})\mid\mathcal{F}_t\bigr]
&\le f(x_t)-\eta_t\|\nabla f(x_t)\|^{2}
+\frac{L}{2}\eta_t^{2}\bigl(\|\nabla f(x_t)\|^{2}+\sigma^{2}\bigr)
\notag\\
&= f(x_t)-\Bigl(\eta_t-\frac{L}{2}\eta_t^{2}\Bigr)\|\nabla f(x_t)\|^{2}
+\frac{L}{2}\eta_t^{2}\sigma^{2}.
\tag{8}
\end{align}
Since $0<\eta\le 1/L$, Lemma~\ref{lem:stepbound} gives $\eta_t\le\eta\le 1/L$, hence
\[
\eta_t-\frac{L}{2}\eta_t^{2}\ge \frac{\eta_t}{2}\ge \frac{\eta}{2(1+\beta c_\alpha G)}.
\tag{9}
\]

\paragraph{[Step 4] Apply the PL inequality (Assumption~\ref{ass:pl}).}
From (9) and the PL condition,
\[
\Bigl(\eta_t-\frac{L}{2}\eta_t^{2}\Bigr)\|\nabla f(x_t)\|^{2}
\ge \frac{\eta_t\mu}{2}\bigl(f(x_t)-f^\star\bigr)
\ge \frac{\eta\mu}{2(1+\beta c_\alpha G)}\bigl(f(x_t)-f^\star\bigr).
\tag{10}
\]

\paragraph{[Step 5] Combine (8) and (10).}
\begin{align}
\mathbb{E}\!\bigl[f(x_{t+1})-f^\star\mid\mathcal{F}_t\bigr]
&\le \Bigl(1-\frac{\eta\mu}{2(1+\beta c_\alpha G)}\Bigr)\bigl(f(x_t)-f^\star\bigr)
+\frac{L}{2}\eta^{2}\sigma^{2}.
\tag{11}
\end{align}

\paragraph{[Step 6] Unroll the recursion.}
Define $\rho:=1-\frac{\eta\mu}{2(1+\beta c_\alpha G)}\in(0,1)$.  Taking total expectation in (11) yields
\[
\mathbb{E}\bigl[f(x_{t+1})-f^\star\bigr]\le \rho\,\mathbb{E}\bigl[f(x_t)-f^\star\bigr]+\frac{L}{2}\eta^{2}\sigma^{2}.
\]
Iterating this inequality $T$ times gives
\[
\mathbb{E}\bigl[f(x_T)-f^\star\bigr]\le \rho^{T}\bigl(f(x_0)-f^\star\bigr)
+\frac{L}{2}\eta^{2}\sigma^{2}\sum_{k=0}^{T-1}\rho^{k}.
\tag{12}
\]

\paragraph{[Step 7] Sum the geometric series.}
Since $\sum_{k=0}^{T-1}\rho^{k}\le \frac{1}{1-\rho}= \frac{2(1+\beta c_\alpha G)}{\eta\mu}$,
\[
\frac{L}{2}\eta^{2}\sigma^{2}\sum_{k=0}^{T-1}\rho^{k}
\le \frac{L}{2}\eta^{2}\sigma^{2}\cdot\frac{2(1+\beta c_\alpha G)}{\eta\mu}
= \frac{L\eta(1+\beta c_\alpha G)}{\mu}\sigma^{2}.
\tag{13}
\]

\paragraph{[Step 8] Choose $\eta$ small enough so that $L\eta(1+\beta c_\alpha G)\le\frac{1}{2}$.}
The condition $\eta\le \frac{1}{2L(1+\beta c_\alpha G)}$ guarantees
\[
\frac{L\eta(1+\beta c_\alpha G)}{\mu}\sigma^{2}\le \frac{\eta\sigma^{2}}{2\mu}.
\tag{14}
\]

\paragraph{[Step 9] Substitute (14) into (12).}
\[
\mathbb{E}\bigl[f(x_T)-f^\star\bigr]\le
\rho^{T}\bigl(f(x_0)-f^\star\bigr)+\frac{\eta\sigma^{2}}{2\mu}.
\]
Recall $\rho=1-\frac{\eta\mu}{2(1+\beta c_\alpha G)}\le 1-\frac{\eta\mu}{2}$ because $1+\beta c_\alpha G\ge1$.  Hence
\[
\rho^{T}\le \bigl(1-\tfrac{\eta\mu}{2}\bigr)^{T}.
\]
This completes the proof of \eqref{5}.  \hfill $\square$

\section*{Rate Derivation (With Constants)}
From Theorem~\ref{thm:linear} we obtain two regimes:

\begin{itemize}
\item \textbf{Deterministic (full‑batch) case $\sigma=0$.} 
\[
\mathbb{E}[f(x_T)-f^\star]\le\bigl(1-\tfrac{\eta\mu}{2}\bigr)^{T}\bigl(f(x_0)-f^\star\bigr),
\]
i.e. linear convergence with factor $1-\eta\mu/2$.

\item \textbf{Stochastic case.} After $T\ge\mathcal{O}\!\bigl(\frac{1}{\eta\mu}\log\frac{1}{\epsilon}\bigr)$ iterations,
\[
\mathbb{E}[f(x_T)-f^\star]\le\epsilon+\frac{\eta\sigma^{2}}{2\mu}.
\]
Choosing $\eta=\Theta\bigl(\frac{\epsilon\mu}{\sigma^{2}}\bigr)$ balances the two terms and yields an overall complexity
\[
\mathcal{O}\!\Bigl(\frac{L}{\mu}\log\frac{1}{\epsilon}+\frac{\sigma^{2}}{\mu\epsilon}\Bigr).
\]
When $\tilde g_t$ is computed by a SARAH/SPIDER inner loop, $\sigma^{2}=O(\epsilon^{2})$ and the second term vanishes, recovering pure linear convergence.
\end{itemize}

\section*{Special Cases}
\begin{enumerate}
\item \textbf{No adaptivity ($\beta=0$).} The algorithm reduces to SGD with EMA momentum; the convergence bound coincides with classical PL results.
\item \textbf{Full momentum ($\alpha=1$).} Then $v_t=\tilde g_t$ and $\hat v_t=\tilde g_t$; the stepsize becomes $\eta/(1+\beta\|\tilde g_t\|)$, a well‑known gradient‑norm adaptive rule.
\item \textbf{Variance‑reduced inner loop.} If $\tilde g_t$ is a SARAH estimator satisfying $\mathbb{E}[\|\tilde g_t-\nabla f(x_t)\|^{2}]\le\frac{L^{2}}{b}\|x_t-x_{t-1}\|^{2}$, the variance term $\sigma^{2}$ decays proportionally to $\|x_t-x_{t-1}\|^{2}$, leading to a tighter bound than \eqref{5}.
\end{enumerate}

\bigskip
\noindent\textbf{References}\\
\begin{thebibliography}{9}
\bibitem{karimi2016linear}
H.~Karimi, J.~Nutini, and M.~Schmidt,
``Linear convergence of gradient and proximal‑gradient methods under the
Polyak–Łojasiewicz condition,''
\emph{J. Optim. Theory Appl.}, vol.~162, no.~2, pp.~307–331, 2016.
\end{thebibliography}

\end{document}